\documentclass[12pt,a4paper]{article}
\usepackage[utf8]{inputenc}
\usepackage[turkish,english]{babel}
\usepackage{amsmath,amssymb,amsfonts}
\usepackage{graphicx}
\usepackage{hyperref}
\usepackage{listings}
\usepackage{xcolor}
\usepackage{algorithm}
\usepackage{algorithmic}
\usepackage{geometry}
\geometry{left=2.5cm,right=2.5cm,top=2.5cm,bottom=2.5cm}

\lstset{
    basicstyle=\ttfamily\small,
    keywordstyle=\color{blue},
    commentstyle=\color{gray},
    stringstyle=\color{red},
    breaklines=true,
    showstringspaces=false
}

\title{\textbf{Cross-View Geo-Localization with RGBD:\\Technical Documentation and Theoretical Background}}
\author{University-1652 Baseline with MiDaS Depth Estimation}
\date{\today}

\begin{document}

\maketitle
\tableofcontents
\newpage

%===================================================================
\section{Introduction}
%===================================================================

Bu dokümantasyon, satellite-to-drone cross-view geo-localization sistemi için geliştirilen RGBD (RGB + Depth) tabanlı deep learning pipeline'ının detaylı teknik açıklamasını içermektedir.

\subsection{Problem Tanımı}
\textbf{Cross-view geo-localization}, farklı bakış açılarından (viewpoints) alınmış görüntüler arasında eşleştirme yaparak lokasyon tespiti problemidir.

\begin{itemize}
    \item \textbf{Query}: Satellite (uydu) görüntüsü
    \item \textbf{Gallery}: Drone (İHA) görüntüsü
    \item \textbf{Amaç}: Query görüntüye en benzer gallery görüntüyü bul
\end{itemize}

\subsection{Araştırma Konuları (Research Keywords)}
\begin{itemize}
    \item Cross-view image matching
    \item Geo-localization from aerial images
    \item Multi-view metric learning
    \item Siamese networks
    \item Circle Loss
    \item Monocular depth estimation
    \item RGBD deep learning
    \item Transfer learning with pretrained CNNs
    \item Random erasing augmentation
    \item Domain adaptation
\end{itemize}

%===================================================================
\section{Model Architecture (model.py, model\_rgbd.py)}
%===================================================================

\subsection{GeM Pooling}

\textbf{Generalized Mean (GeM) Pooling}, feature map'lerde spatial pooling için kullanılan gelişmiş bir yöntemdir.

\subsubsection{Matematiksel Tanım}
\begin{equation}
    \text{GeM}(X) = \left(\frac{1}{|X|} \sum_{x \in X} x^p\right)^{1/p}
\end{equation}

Burada:
\begin{itemize}
    \item $X$: Feature map ($(H \times W)$ boyutunda)
    \item $p$: Öğrenilebilir parametre (learnable power parameter)
    \item $|X|$: Feature map'teki eleman sayısı
\end{itemize}

\subsubsection{Özel Durumlar}
\begin{align}
    p = 1 &\implies \text{Average Pooling} \\
    p = \infty &\implies \text{Max Pooling} \\
    p = 3 &\implies \text{Default GeM} \text{ (tipik başlangıç değeri)}
\end{align}

\subsubsection{PyTorch Implementasyonu}
\begin{lstlisting}[language=Python]
class GeM(nn.Module):
    def __init__(self, dim=2048, p=3, eps=1e-6):
        super(GeM, self).__init__()
        # p parametresi gradient descent ile öğrenilir
        self.p = nn.Parameter(torch.ones(dim)*p, requires_grad=True)
        self.eps = eps
        
    def forward(self, x):
        # x shape: (B, C, H, W)
        x = x.clamp(min=self.eps).pow(self.p)  # x^p
        x = F.avg_pool2d(x, (x.size(-2), x.size(-1)))  # Global pooling
        x = x.view(x.size(0), x.size(1))  # (B, C)
        x = x.pow(1./self.p)  # p-th root
        return x
\end{lstlisting}

\subsubsection{Avantajları}
\begin{itemize}
    \item \textbf{Adaptif}: $p$ parametresi eğitim sırasında öğrenilir
    \item \textbf{Diskriminatif}: Average ve max pooling arasında optimal dengeyi bulur
    \item \textbf{Image retrieval}: Retrieval tasklerinde SOTA performans
\end{itemize}

\textbf{Teorik Kaynak}: Radenović et al., "Fine-tuning CNN Image Retrieval with No Human Annotation", TPAMI 2019

%-------------------------------------------------------------------
\subsection{ClassBlock: Feature Embedding Layer}

\textbf{ClassBlock}, feature extraction sonrası classification head'i oluşturur.

\subsubsection{Mimari Yapı}
\begin{equation}
    \text{ClassBlock}: \mathbb{R}^{2048} \xrightarrow{\text{Linear}} \mathbb{R}^{512} \xrightarrow{\text{BN}} \xrightarrow{\text{Dropout}} \xrightarrow{\text{Linear}} \mathbb{R}^{C}
\end{equation}

\begin{itemize}
    \item \textbf{Input}: Feature vector $f \in \mathbb{R}^{2048}$ (ResNet50 output)
    \item \textbf{Bottleneck}: $f' \in \mathbb{R}^{512}$ (embedding space)
    \item \textbf{Output}: Logits $y \in \mathbb{R}^{C}$ ($C$ = number of classes)
\end{itemize}

\subsubsection{Bileşenler}
\begin{enumerate}
    \item \textbf{Linear (2048 → 512)}: Dimensionality reduction
    \item \textbf{BatchNorm1d}: Normalization for stable training
    \item \textbf{LeakyReLU (optional)}: Non-linearity, $\alpha = 0.1$
    \item \textbf{Dropout}: Regularization, $p = 0.5$
    \item \textbf{Classifier Linear (512 → C)}: Final classification layer
\end{enumerate}

\subsubsection{Weight Initialization}
\begin{itemize}
    \item \textbf{Kaiming Initialization} (He initialization):
    \begin{equation}
        W \sim \mathcal{N}\left(0, \sqrt{\frac{2}{n_{\text{in}}}}\right)
    \end{equation}
    
    \item \textbf{Classifier Initialization}:
    \begin{equation}
        W \sim \mathcal{N}(0, 0.001^2), \quad b = 0
    \end{equation}
\end{itemize}

\textbf{Teorik Kaynak}: He et al., "Delving Deep into Rectifiers", ICCV 2015

%-------------------------------------------------------------------
\subsection{RGBD Model Architecture (model\_rgbd.py)}

\subsubsection{4-Channel Input Conversion}

Standard ResNet50 3-channel (RGB) inputu alır. RGBD için 4. kanalı (depth) eklemek gerekir.

\textbf{Fonksiyon}: \texttt{convert\_conv1\_to\_4channel}

\paragraph{Adımlar:}
\begin{enumerate}
    \item Orijinal \texttt{conv1} ağırlıkları: $W_{\text{orig}} \in \mathbb{R}^{64 \times 3 \times 7 \times 7}$
    \item Yeni \texttt{conv1} oluştur: $W_{\text{new}} \in \mathbb{R}^{64 \times 4 \times 7 \times 7}$
    \item RGB kanallarını kopyala:
    \begin{equation}
        W_{\text{new}}[:, 0:3, :, :] = W_{\text{orig}}
    \end{equation}
    
    \item Depth kanalını initialize et (RGB ortalaması):
    \begin{equation}
        W_{\text{new}}[:, 3, :, :] = \frac{1}{3}\sum_{c=0}^{2} W_{\text{orig}}[:, c, :, :]
    \end{equation}
\end{enumerate}

\subsubsection{Neden RGB Ortalaması?}
\begin{itemize}
    \item \textbf{Transfer learning}: Depth kanalı sıfırdan başlamaz
    \item \textbf{Benzer semantik}: Depth ve RGB intensity korelasyonlu
    \item \textbf{Hızlı convergence}: Random initialization'dan daha iyi
\end{itemize}

%-------------------------------------------------------------------
\subsection{Two-View Network}

\textbf{two\_view\_net\_rgbd} sınıfı, satellite (RGBD) ve drone (RGB) için ayrı encoder kullanır.

\subsubsection{Matematiksel Formülasyon}
\begin{align}
    x_{\text{sat}} &\in \mathbb{R}^{B \times 4 \times H \times W} \quad \text{(Satellite RGBD)} \\
    x_{\text{drone}} &\in \mathbb{R}^{B \times 3 \times H \times W} \quad \text{(Drone RGB)}
\end{align}

\begin{align}
    f_{\text{sat}} &= \phi_1(x_{\text{sat}}) \quad \text{(ResNet50-RGBD encoder)} \\
    f_{\text{drone}} &= \phi_2(x_{\text{drone}}) \quad \text{(ResNet50-RGB encoder)}
\end{align}

Her iki feature $f_{\text{sat}}, f_{\text{drone}} \in \mathbb{R}^{B \times 2048}$ aynı classifier'a girer:
\begin{equation}
    y = \psi(f), \quad \psi: \mathbb{R}^{2048} \to \mathbb{R}^{512} \to \mathbb{R}^{C}
\end{equation}

\subsubsection{Shared Classifier}
\begin{itemize}
    \item \textbf{Embedding space paylaşılır}: Her iki view aynı metric space'e map edilir
    \item \textbf{Cross-view matching}: $\text{sim}(f_{\text{sat}}, f_{\text{drone}}) = f_{\text{sat}}^T f_{\text{drone}}$
    \item \textbf{Circle Loss}: Embedding space'de intra-class compact, inter-class separable
\end{itemize}

%===================================================================
\section{Dataset Loading (dataset\_rgbd.py)}
%===================================================================

\subsection{RGBDSatelliteDataset Sınıfı}

\subsubsection{Veri Formatı}
\begin{itemize}
    \item \textbf{RGB}: \texttt{satellite/0000/image001.jpg} $(H \times W \times 3)$
    \item \textbf{Depth}: \texttt{satellite\_depth/0000/image001\_depth.jpg} $(H \times W \times 1)$
\end{itemize}

\subsubsection{Pipeline}
\begin{algorithm}[H]
\caption{RGBD Data Loading}
\begin{algorithmic}[1]
\STATE Load RGB image: $I_{\text{rgb}} \in [0, 255]^{H \times W \times 3}$
\STATE Load Depth image: $I_{\text{depth}} \in [0, 255]^{H \times W}$ (grayscale)
\STATE Apply transform to RGB: $T_{\text{rgb}}(I_{\text{rgb}}) \to \mathbb{R}^{3 \times 256 \times 256}$
\STATE Apply transform to Depth: $T_{\text{depth}}(I_{\text{depth}}) \to \mathbb{R}^{1 \times 256 \times 256}$
\STATE Concatenate: $I_{\text{rgbd}} = [I_{\text{rgb}}, I_{\text{depth}}] \in \mathbb{R}^{4 \times 256 \times 256}$
\RETURN $I_{\text{rgbd}}$, label
\end{algorithmic}
\end{algorithm}

\subsubsection{Transform Pipeline}
\begin{lstlisting}[language=Python]
transform = transforms.Compose([
    transforms.Resize((256, 256)),
    transforms.RandomHorizontalFlip(),
    transforms.RandomRotation(90),  # Satellite için
    transforms.ToTensor(),
    transforms.Normalize(mean=[0.485, 0.456, 0.406], 
                         std=[0.229, 0.224, 0.225])
])
\end{lstlisting}

\textbf{Not}: Depth kanalı için normalization RGB ile aynı (ImageNet statistics).

%===================================================================
\section{Loss Functions}
%===================================================================

\subsection{Circle Loss (circle\_loss.py)}

\textbf{Circle Loss}, metric learning için tasarlanmış gelişmiş bir loss fonksiyonudur.

\subsubsection{Motivasyon}
Traditional triplet loss ve softmax loss'un sınırlamalarını aşmak:
\begin{itemize}
    \item \textbf{Triplet loss}: Mining stratejisi gerektirir (hard negative)
    \item \textbf{Softmax}: Embedding optimization için suboptimal
    \item \textbf{Circle loss}: Her iki avantajı birleştirir
\end{itemize}

\subsubsection{Matematiksel Formülasyon}

Normalized feature space'de:
\begin{equation}
    \mathbf{f}_i = \frac{\mathbf{z}_i}{\|\mathbf{z}_i\|_2}, \quad \mathbf{f}_i \in \mathbb{R}^{512}
\end{equation}

\paragraph{Similarity Matrix:}
\begin{equation}
    S_{ij} = \mathbf{f}_i^T \mathbf{f}_j \in [-1, 1]
\end{equation}

\paragraph{Positive ve Negative Pairs:}
\begin{align}
    \mathcal{P}_i &= \{j \mid y_i = y_j, j \neq i\} \quad \text{(Positive pairs)} \\
    \mathcal{N}_i &= \{j \mid y_i \neq y_j\} \quad \text{(Negative pairs)}
\end{align}

\paragraph{Circle Loss Formülü:}
\begin{equation}
    \mathcal{L}_{\text{circle}} = \log\left[1 + \sum_{n \in \mathcal{N}_i} \exp(\gamma \alpha_n (s_n - \Delta_n)) \cdot \sum_{p \in \mathcal{P}_i} \exp(-\gamma \alpha_p (s_p - \Delta_p))\right]
\end{equation}

Burada:
\begin{align}
    \alpha_n &= \max(0, s_n + m) \quad \text{(Negative weight)} \\
    \alpha_p &= \max(0, -s_p + 1 + m) \quad \text{(Positive weight)} \\
    \Delta_n &= m, \quad \Delta_p = 1 - m \\
    m &= 0.25 \quad \text{(Margin parameter)} \\
    \gamma &= 256 \quad \text{(Scale factor)}
\end{align}

\subsubsection{Intuition}
\begin{itemize}
    \item \textbf{Positive pairs}: $s_p \to 1$ (maximize similarity)
    \item \textbf{Negative pairs}: $s_n \to -1$ (minimize similarity)
    \item \textbf{Adaptive weights}: Zor örneklere daha fazla ağırlık
    \item \textbf{Margin $m$}: Decision boundary için buffer
\end{itemize}

\textbf{Teorik Kaynak}: Sun et al., "Circle Loss: A Unified Perspective of Pair Similarity Optimization", CVPR 2020

%-------------------------------------------------------------------
\subsection{Cross-Entropy Loss}

Standard classification loss:
\begin{equation}
    \mathcal{L}_{\text{CE}} = -\frac{1}{B}\sum_{i=1}^{B} \log\frac{\exp(z_{y_i})}{\sum_{j=1}^{C} \exp(z_j)}
\end{equation}

Burada:
\begin{itemize}
    \item $z \in \mathbb{R}^{C}$: Logits (classifier output)
    \item $y_i$: Ground truth label
    \item $B$: Batch size
    \item $C$: Number of classes (701 for University-1652)
\end{itemize}

%-------------------------------------------------------------------
\subsection{Combined Loss}

Training sırasında her iki loss birleştirilir:
\begin{equation}
    \mathcal{L}_{\text{total}} = \mathcal{L}_{\text{CE}} + \lambda \mathcal{L}_{\text{circle}}
\end{equation}

Tipik olarak $\lambda = 1$ (equal weighting).

%===================================================================
\section{Data Augmentation}
%===================================================================

\subsection{Random Erasing (random\_erasing.py)}

\textbf{Random Erasing}, overfitting'i azaltmak için görüntünün rastgele bir bölgesini maskeler.

\subsubsection{Algoritma}
\begin{algorithm}[H]
\caption{Random Erasing}
\begin{algorithmic}[1]
\STATE \textbf{Input}: Image $I \in \mathbb{R}^{C \times H \times W}$, probability $p$
\STATE Generate random number $r \sim \mathcal{U}(0,1)$
\IF{$r > p$}
    \RETURN $I$ \quad (no erasing)
\ENDIF
\STATE Select area ratio: $S_e \sim \mathcal{U}(s_l, s_h)$ where $s_l=0.02, s_h=0.4$
\STATE Select aspect ratio: $r_e \sim \mathcal{U}(r_1, 1/r_1)$ where $r_1=0.3$
\STATE Compute dimensions: $h_e = \sqrt{S_e \cdot H \cdot W \cdot r_e}$, $w_e = \sqrt{S_e \cdot H \cdot W / r_e}$
\STATE Random position: $(x_1, y_1) \sim \mathcal{U}(0, H-h_e) \times \mathcal{U}(0, W-w_e)$
\STATE Erase region: $I[:, x_1:x_1+h_e, y_1:y_1+w_e] = \mu$ (mean color)
\RETURN $I$
\end{algorithmic}
\end{algorithm}

\subsubsection{Parametreler}
\begin{itemize}
    \item $p = 0.5$: Erasing probability
    \item $s_l = 0.02, s_h = 0.4$: Area ratio bounds
    \item $r_1 = 0.3$: Aspect ratio bound
    \item $\mu = [0.4914, 0.4822, 0.4465]$: ImageNet mean
\end{itemize}

\textbf{Teorik Kaynak}: Zhong et al., "Random Erasing Data Augmentation", AAAI 2020

%-------------------------------------------------------------------
\subsection{Color Jitter}

\begin{equation}
    I' = \alpha \cdot I + \beta
\end{equation}

\begin{itemize}
    \item \textbf{Brightness}: $\alpha \sim \mathcal{U}(0.9, 1.1)$
    \item \textbf{Contrast}: Scaling factor $\sim \mathcal{U}(0.9, 1.1)$
    \item \textbf{Saturation}: HSV space'de $\pm 10\%$
    \item \textbf{Hue}: $\pm 0$ (tipik olarak devre dışı)
\end{itemize}

%-------------------------------------------------------------------
\subsection{Geometric Augmentations}

\subsubsection{Satellite-Specific}
\begin{itemize}
    \item \textbf{RandomRotation(90°)}: Uydu görüntüsü yön invaryant
    \item \textbf{RandomAffine}: Rotation, translation, scale
    \item \textbf{RandomHorizontalFlip}: Mirror symmetry
\end{itemize}

\subsubsection{Drone-Specific}
\begin{itemize}
    \item \textbf{RandomCrop}: Simüle edilmiş zoom
    \item \textbf{RandomHorizontalFlip}: View variation
    \item \textbf{Padding + Crop}: Edge artifacts'ı önle
\end{itemize}

%===================================================================
\section{Training Pipeline (train\_cvusa.py)}
%===================================================================

\subsection{Hyperparameters}

\begin{table}[h]
\centering
\begin{tabular}{|l|c|l|}
\hline
\textbf{Parameter} & \textbf{Value} & \textbf{Description} \\
\hline
Batch Size & 8 & Per-GPU batch size \\
Learning Rate & 0.01 & Initial LR \\
Optimizer & SGD & Momentum = 0.9 \\
Weight Decay & $5 \times 10^{-4}$ & L2 regularization \\
Epochs & 60 & Total training epochs \\
Scheduler & StepLR & Decay at [40, 50] \\
Decay Factor & 0.1 & LR multiplier \\
Dropout & 0.5 & Classifier dropout \\
Image Size & $384 \times 384$ & Input resolution \\
\hline
\end{tabular}
\caption{Training Hyperparameters}
\end{table}

%-------------------------------------------------------------------
\subsection{Learning Rate Schedule}

\begin{equation}
    \text{LR}(\text{epoch}) = \begin{cases}
        0.01 & \text{epoch} \in [0, 40) \\
        0.001 & \text{epoch} \in [40, 50) \\
        0.0001 & \text{epoch} \in [50, 60]
    \end{cases}
\end{equation}

\textbf{Warm-up} (optional):
\begin{equation}
    \text{LR}(\text{epoch}) = \frac{\text{epoch}}{K} \cdot \text{LR}_{\text{base}}, \quad \text{epoch} < K
\end{equation}
Tipik $K = 5$ (5 epoch warm-up).

%-------------------------------------------------------------------
\subsection{Training Loop}

\begin{algorithm}[H]
\caption{Two-View Training}
\begin{algorithmic}[1]
\FOR{epoch = 1 to 60}
    \FOR{each batch in dataloader}
        \STATE Sample satellite batch: $X_{\text{sat}}, Y_{\text{sat}}$
        \STATE Sample drone batch: $X_{\text{drone}}, Y_{\text{drone}}$
        \STATE Forward pass: $Z_{\text{sat}}, F_{\text{sat}} = \text{model}(X_{\text{sat}}, \text{None})$
        \STATE Forward pass: $Z_{\text{drone}}, F_{\text{drone}} = \text{model}(\text{None}, X_{\text{drone}})$
        \STATE Compute CE loss: $\mathcal{L}_{\text{CE,sat}} + \mathcal{L}_{\text{CE,drone}}$
        \STATE Normalize features: $F_{\text{sat}} = F_{\text{sat}} / \|F_{\text{sat}}\|_2$
        \STATE Compute Circle loss: $\mathcal{L}_{\text{circle}}(F_{\text{sat}}, F_{\text{drone}}, Y)$
        \STATE Total loss: $\mathcal{L} = \mathcal{L}_{\text{CE}} + \mathcal{L}_{\text{circle}}$
        \STATE Backward: $\nabla_\theta \mathcal{L}$
        \STATE Optimizer step
    \ENDFOR
    \STATE Update learning rate
    \STATE Save checkpoint
\ENDFOR
\end{algorithmic}
\end{algorithm}

%===================================================================
\section{Evaluation Metrics (evaluate\_gpu.py)}
%===================================================================

\subsection{Cumulative Matching Characteristic (CMC)}

\textbf{CMC curve}, retrieval sistemlerinde kullanılan temel metriklerdir.

\subsubsection{Tanım}
\begin{equation}
    \text{CMC@K} = \frac{1}{N_q} \sum_{i=1}^{N_q} \mathbb{1}[\text{rank}(i) \leq K]
\end{equation}

Burada:
\begin{itemize}
    \item $N_q$: Query sayısı
    \item $\text{rank}(i)$: $i$-inci query için doğru match'in sıralaması
    \item $K$: Rank threshold (1, 5, 10)
\end{itemize}

\subsubsection{Yorumlama}
\begin{itemize}
    \item \textbf{Rank-1}: İlk tahmin doğru mu? (Top-1 accuracy)
    \item \textbf{Rank-5}: İlk 5 tahmin içinde doğru var mı?
    \item \textbf{Rank-10}: İlk 10 tahmin içinde doğru var mı?
\end{itemize}

%-------------------------------------------------------------------
\subsection{Mean Average Precision (mAP)}

\subsubsection{Average Precision (AP)}
Tek bir query için:
\begin{equation}
    \text{AP} = \frac{1}{N_{\text{rel}}} \sum_{k=1}^{N} P(k) \cdot \text{rel}(k)
\end{equation}

\begin{itemize}
    \item $N_{\text{rel}}$: Relevant (doğru) sonuç sayısı
    \item $P(k) = \frac{\text{\# relevant in top-}k}{k}$: Precision at rank $k$
    \item $\text{rel}(k) = 1$ if rank $k$ relevant, else 0
\end{itemize}

\subsubsection{Mean Average Precision}
\begin{equation}
    \text{mAP} = \frac{1}{N_q} \sum_{i=1}^{N_q} \text{AP}_i
\end{equation}

\textbf{Yorumlama}: mAP, tüm querylerin ortalama precision'ını ölçer. Yüksek mAP, sistemin genel retrieval kalitesinin iyi olduğunu gösterir.

%-------------------------------------------------------------------
\subsection{Evaluation Pipeline}

\begin{algorithm}[H]
\caption{Cross-View Evaluation}
\begin{algorithmic}[1]
\STATE Extract query features: $\{f_i^q\}_{i=1}^{N_q}$ from satellite
\STATE Extract gallery features: $\{f_j^g\}_{j=1}^{N_g}$ from drone
\STATE Normalize: $f \leftarrow f / \|f\|_2$
\FOR{each query $i$}
    \STATE Compute similarity: $s_{ij} = (f_i^q)^T f_j^g, \forall j$
    \STATE Sort gallery by similarity (descending)
    \STATE Find rank of correct match
    \STATE Compute AP
\ENDFOR
\STATE Compute CMC@1, CMC@5, CMC@10
\STATE Compute mAP
\end{algorithmic}
\end{algorithm}

%===================================================================
\section{MiDaS Depth Estimation}
%===================================================================

\subsection{Monocular Depth Estimation}

\textbf{MiDaS} (Mixed Data Sampling), tek bir RGB görüntüden depth map tahmin eden bir modeldir.

\subsubsection{Mimari}
\begin{itemize}
    \item \textbf{Encoder}: EfficientNet-based backbone
    \item \textbf{Decoder}: Multi-scale feature fusion
    \item \textbf{Output}: Relative depth map (inverse depth)
\end{itemize}

\subsubsection{Matematiksel Model}
\begin{equation}
    D = \mathcal{F}(I_{\text{RGB}})
\end{equation}

Burada:
\begin{itemize}
    \item $I_{\text{RGB}} \in \mathbb{R}^{H \times W \times 3}$: Input RGB image
    \item $D \in \mathbb{R}^{H \times W}$: Depth map (relative values)
    \item $\mathcal{F}$: MiDaS neural network
\end{itemize}

\subsubsection{Depth Normalization}
MiDaS çıktısı relative depth olduğundan, normalize edilmesi gerekir:
\begin{equation}
    D_{\text{norm}} = 255 \cdot \frac{D - D_{\min}}{D_{\max} - D_{\min}}
\end{equation}

%-------------------------------------------------------------------
\subsection{MiDaS Small Model}

\begin{itemize}
    \item \textbf{Model boyutu}: $\sim$21MB
    \item \textbf{Inference hızı}: $\sim$0.1s/image (GPU)
    \item \textbf{Input size}: Any resolution (internally resized)
    \item \textbf{Output}: Same resolution as input
\end{itemize}

\textbf{Kullanım}:
\begin{lstlisting}[language=Python]
midas = torch.hub.load("intel-isl/MiDaS", "MiDaS_small")
transform = torch.hub.load("intel-isl/MiDaS", "transforms")

img = cv2.imread("satellite.jpg")
input_batch = transform.small_transform(img)
depth = midas(input_batch)
\end{lstlisting}

\textbf{Teorik Kaynak}: Ranftl et al., "Towards Robust Monocular Depth Estimation: Mixing Datasets for Zero-shot Cross-dataset Transfer", TPAMI 2020

%===================================================================
\section{Theoretical Background}
%===================================================================

\subsection{Metric Learning}

\textbf{Amaç}: Feature space'de benzer örnekler yakın, farklı örnekler uzak olacak şekilde metric öğren.

\subsubsection{Distance Metrics}
\begin{itemize}
    \item \textbf{Euclidean}: $d(x, y) = \|x - y\|_2$
    \item \textbf{Cosine}: $d(x, y) = 1 - \frac{x^T y}{\|x\|\|y\|}$
    \item \textbf{Mahalanobis}: $d(x, y) = \sqrt{(x-y)^T M (x-y)}$
\end{itemize}

Bu projede \textbf{cosine similarity} kullanılıyor:
\begin{equation}
    \text{sim}(x, y) = \frac{x^T y}{\|x\|\|y\|} \in [-1, 1]
\end{equation}

%-------------------------------------------------------------------
\subsection{Transfer Learning}

\textbf{ImageNet pretrained weights}, bu projede transfer learning için kullanılır.

\subsubsection{Neden Transfer Learning?}
\begin{itemize}
    \item \textbf{Veri kısıtı}: University-1652 nispeten küçük ($\sim$50K images)
    \item \textbf{Generalization}: ImageNet'te öğrenilen low-level features evrensel
    \item \textbf{Hız}: Random initialization'dan çok daha hızlı converge
\end{itemize}

\subsubsection{Fine-tuning Strategy}
\begin{enumerate}
    \item Load pretrained ResNet50 (ImageNet weights)
    \item Freeze initial layers (conv1-layer2) - optional
    \item Train classifier + final layers (layer3-layer4)
    \item Gradually unfreeze early layers
\end{enumerate}

%-------------------------------------------------------------------
\subsection{Cross-View Matching Challenges}

\begin{itemize}
    \item \textbf{Viewpoint variation}: Satellite nadir view, drone oblique view
    \item \textbf{Scale difference}: Satellite çok daha geniş alan kaplar
    \item \textbf{Appearance gap}: Lighting, weather, season differences
    \item \textbf{Occlusion}: Buildings, trees farklı görünür
\end{itemize}

\textbf{RGBD'nin katkısı}:
\begin{itemize}
    \item Depth, 3D structure bilgisi sağlar
    \item Viewpoint invariant features öğrenmeyi kolaylaştırır
    \item Buildings height, terrain topology gibi geometric cues
\end{itemize}

%===================================================================
\section{Advanced Research Topics}
%===================================================================

\subsection{Future Directions}

\subsubsection{Attention Mechanisms}
\begin{itemize}
    \item \textbf{Spatial Attention}: Önemli bölgelere focus
    \item \textbf{Channel Attention}: Informative channels'ı vurgula
    \item \textbf{Cross-Attention}: Satellite-drone feature alignment
\end{itemize}

\subsubsection{Semi-Supervised Learning}
\begin{itemize}
    \item \textbf{Pseudo-labeling}: Unlabeled data'dan yararlan
    \item \textbf{Consistency regularization}: Augmented views arasında consistency
    \item \textbf{Contrastive learning}: SimCLR, MoCo benzeri approaches
\end{itemize}

\subsubsection{Multi-Modal Fusion}
\begin{itemize}
    \item \textbf{Early fusion}: RGBD gibi, input seviyesinde birleştir
    \item \textbf{Late fusion}: Feature seviyesinde birleştir
    \item \textbf{Attention fusion}: Learned weights ile modal combination
\end{itemize}

%===================================================================
\section{Araştırma Kaynakları (Research Keywords)}
%===================================================================

\subsection{Temel Konular}
\begin{itemize}
    \item \textbf{Cross-view image matching}
    \item \textbf{Geo-localization from aerial images}
    \item \textbf{Multi-view metric learning}
    \item \textbf{Siamese networks}
    \item \textbf{Triplet loss and its variants}
    \item \textbf{Circle loss}
    \item \textbf{Monocular depth estimation}
    \item \textbf{RGBD deep learning}
    \item \textbf{Transfer learning with CNNs}
\end{itemize}

\subsection{İleri Seviye}
\begin{itemize}
    \item \textbf{Domain adaptation for cross-view matching}
    \item \textbf{Attention mechanisms in computer vision}
    \item \textbf{Self-supervised learning}
    \item \textbf{Contrastive learning}
    \item \textbf{Image retrieval methods}
    \item \textbf{Visual place recognition}
    \item \textbf{Feature aggregation techniques}
    \item \textbf{Hard negative mining strategies}
\end{itemize}

\subsection{Dataset ve Benchmarklar}
\begin{itemize}
    \item \textbf{University-1652 dataset}
    \item \textbf{CVUSA dataset}
    \item \textbf{CVACT dataset}
    \item \textbf{VIGOR dataset}
\end{itemize}

%===================================================================
\section{Pratik Uygulama Önerileri}
%===================================================================

\subsection{Hiperparametre Tuning}

\begin{table}[h]
\centering
\begin{tabular}{|l|c|c|}
\hline
\textbf{Parameter} & \textbf{Default} & \textbf{Tuning Range} \\
\hline
Learning Rate & 0.01 & [0.001, 0.1] \\
Batch Size & 8 & [4, 16, 32] \\
Dropout & 0.5 & [0.3, 0.7] \\
Weight Decay & $5 \times 10^{-4}$ & [$10^{-5}$, $10^{-3}$] \\
Circle Loss $m$ & 0.25 & [0.1, 0.5] \\
Circle Loss $\gamma$ & 256 & [64, 512] \\
\hline
\end{tabular}
\caption{Hyperparameter Tuning Ranges}
\end{table}

\subsection{Debugging Checklist}
\begin{enumerate}
    \item \textbf{Data loading}: Print batch shapes, visualize samples
    \item \textbf{Model output}: Check logits range, NaN values
    \item \textbf{Gradients}: Monitor gradient norms (exploding/vanishing)
    \item \textbf{Loss}: Plot training loss curve
    \item \textbf{Accuracy}: Compute train/val accuracy each epoch
    \item \textbf{Features}: Visualize t-SNE of learned embeddings
\end{enumerate}

%===================================================================
\section{Sonuç}
%===================================================================

Bu dokümantasyon, satellite-to-drone cross-view geo-localization için RGBD-based deep learning sisteminin detaylı teknik açıklamasını sunmaktadır. Model mimarisi, loss fonksiyonları, data augmentation, training pipeline ve evaluation metrikleri teorik temelleriyle birlikte açıklanmıştır.

\subsection{Ana Katkılar}
\begin{enumerate}
    \item \textbf{4-channel RGBD input}: MiDaS depth estimation ile 3D bilgi eklenmesi
    \item \textbf{Two-view architecture}: Satellite ve drone için ayrı encoders
    \item \textbf{Circle Loss}: Metric learning için state-of-the-art loss
    \item \textbf{Transfer learning}: ImageNet pretrained weights kullanımı
\end{enumerate}

\subsection{Beklenen Sonuçlar}
\begin{itemize}
    \item \textbf{Rank-1 Accuracy}: 35-40\%
    \item \textbf{Rank-5 Accuracy}: 60-65\%
    \item \textbf{mAP}: 43-48\%
\end{itemize}

\vspace{1cm}
\noindent
\textit{Bu doküman, University-1652 Baseline projesinin teknik implementasyonunu ve teorik altyapısını detaylı olarak açıklamaktadır. İleri okuma için referans listesindeki makalelere başvurulabilir.}

%===================================================================
\section*{Referanslar}
%===================================================================

\begin{enumerate}
    \item Zheng et al., "University-1652: A Multi-view Multi-source Benchmark for Drone-based Geo-localization", ACM MM 2020
    
    \item Sun et al., "Circle Loss: A Unified Perspective of Pair Similarity Optimization", CVPR 2020
    
    \item Ranftl et al., "Towards Robust Monocular Depth Estimation: Mixing Datasets for Zero-shot Cross-dataset Transfer", TPAMI 2020
    
    \item Radenović et al., "Fine-tuning CNN Image Retrieval with No Human Annotation", TPAMI 2019
    
    \item He et al., "Deep Residual Learning for Image Recognition", CVPR 2016
    
    \item He et al., "Delving Deep into Rectifiers: Surpassing Human-Level Performance on ImageNet Classification", ICCV 2015
    
    \item Zhong et al., "Random Erasing Data Augmentation", AAAI 2020
    
    \item Schroff et al., "FaceNet: A Unified Embedding for Face Recognition and Clustering", CVPR 2015
    
    \item Hermans et al., "In Defense of the Triplet Loss for Person Re-Identification", arXiv 2017
    
    \item Wang et al., "Multi-Similarity Loss with General Pair Weighting for Deep Metric Learning", CVPR 2019
\end{enumerate}

\end{document}
